\documentclass[a4paper,11pt]{letter}
\usepackage[default]{sourcesanspro}
\usepackage{graphicx}
\usepackage{geometry}
\usepackage{fancyhdr}
\usepackage{fontawesome}
\usepackage[utf8]{inputenc}
\usepackage{setspace}
\usepackage[hidelinks]{hyperref}

% Set the page margins
\geometry{a4paper, top=3.3cm, bottom=2.2cm, left=2.5cm, right=2.5cm}

% Setup the header
\pagestyle{fancy}
\fancyhf{} % clear all header fields
\fancyhead[L]{
    {\large \textbf{Prof. Julio Molleda}}\\
    \href{mailto:jmolleda@uniovi.es} jmolleda@uniovi.es\\
    University of Oviedo, Dept. Comp. Sci. and Eng., Gijón, Spain \\

    \medskip
    \textit{Engineering Applications of Artificial Intelligence}
    }
\fancyhead[R]{\textbf{\today}}
\medskip\vspace*{30pt}

%%%%%%  LETTER STARTS HERE  %%%%%%%%%%%%%%%%%%%%%%%%%%%%

\begin{document}

%Salutation
\normalsize Dear Editors,

% Letter body
\small \noindent
We are pleased to submit our manuscript entitled \textbf{``A Shapley-value framework for quantifying module contributions in context-augmented decision support: application to parcel-level land management''} for consideration for publication in \textit{Engineering Applications of Artificial Intelligence}.

\vspace{1pt}
Following the guide for authors, we state the two contributions explicitly. \textbf{The contribution in artificial intelligence} is a framework that models a modular, context-augmented architecture as a cooperative game. It applies Shapley values to attribute response quality to individual data modules. This framework combines a full factorial ablation over module coalitions, a hybrid human-in-the-loop evaluation panel, and non-parametric bootstrap confidence intervals. \textbf{The application in engineering} is a geospatial decision support system for parcel-level land management. It fuses cadastral records, satellite multispectral indices, aerial orthophotos and retrieval-augmented document search to answer natural-language queries about agricultural plots.

\vspace{1pt}
The work addresses a gap that affects any modular system built on augmented context. Although these architectures are increasingly deployed as decision support tools, there is no established method for determining the contribution of each information source to the quality of the output. Our evaluation panel combines two human evaluators and two distinct large language model judges, scoring 143 queries over 15 parcels across five metrics. The orthophoto module emerges as the dominant contributor across that dataset, adding 0.9 out of 5 points (18$\%$) to the mean score. Meanwhile, the external context modules together raise response utility by 32$\%$. The retrieval module exhibits high query-dependent volatility and frequently negative contributions. This indicates that retrieval should be invoked selectively rather than by default. Complementary analyses of context length and inter-judge agreement demonstrate the system resilience to information overload and quantify the divergence between human and machine judges.

\vspace{1pt}
The results were validated using a publicly released dataset. The complete evaluation corpus, containing the queries, the parcel orthophotos, the retrieved multi-domain context, and the generated responses, is deposited in a public repository under the Apache~2.0 license. This allows the reported evaluation to be reproduced independently and the attribution framework to be reused on other modular systems. The Data Availability Statement in the manuscript provides an anonymous reviewer link to the deposit. The permanent identifier of the dataset for the editorial office is \href{https://doi.org/10.57967/hf/10594}{10.57967/hf/10594}, which resolves to the full repository including author metadata. This link is withheld from the anonymized manuscript because it discloses author identity, and will be restored in the published version.

\vspace{1pt}
This manuscript is a substantially extended version of our conference paper presented at the \textit{2025 4th International Conference on Geographic Information and Remote Sensing Technology (GIRST)}\\
(DOI: \href{http://dx.doi.org/10.1109/GIRST67753.2025.11382141}{10.1109/GIRST67753.2025.11382141}), which proposed the multimodal system evaluated here. This journal version shifts the focus from the technical implementation of that system to a mathematically grounded evaluation of its outputs. It also contributes both the attribution methodology and the dataset built for the study. We note that the two papers share no text beyond bibliographic references and the authorship statement.

% \vspace{1pt}
% In accordance with the journal's double anonymized review policy, the manuscript has been anonymized and the title page, carrying author details, acknowledgements and the CRediT statement, is submitted as a separate file.

\vspace{1pt}
We confirm that this manuscript is original, has not been published elsewhere, and is not currently under consideration by any other journal. All authors have approved the manuscript and agree with its submission. The authors declare no competing financial interests or personal relationships that could have appeared to influence the work reported in this paper.

We thank you for considering our work and look forward to your editorial decision.

\bigskip
\noindent
Sincerely,\\
Prof. J. Molleda, on behalf of all authors.\\

\end{document}
